\Needspace{20\baselineskip}\SecondArticle{13}
The proof of this theorem would indeed be very easy if we were allowed to suppose that $Y$ could be resolved into linear factors
\[
(x-A)(x-B)(x-C)(x-D)\, .\, .\, .\, .\, .
\]
For it would then also be certain that $Z$ is the product of all $u-(A+B)x+AB$, and that the determinant of the function $Z$ is the product of the differences between pairs of the quantities
\[
\begin{gathered}
(A+B)x-AB\\
(A+C)x-AC\\
(A+D)x-AD\\
\text{etc.}\\
(B+C)x-BC\\
(B+D)x-BD\\
\text{etc.}\\
(C+D)x-CD\\
\text{etc. etc}
\end{gathered}
\]
But this product cannot vanish identically unless one of its factors itself becomes identically $=0$. This would imply that two of the quantities $A,B,C$ etc. are equal, and thus that the determinant $P$ of the polynomial $Y$ becomes $=0$, contrary to the hypothesis.

\SourceLeaf{055}
Setting aside this argument, which clearly starts by begging the question just as in Article 6, we now proceed directly to give a sound proof of the theorem of Article 12.

\SecondArticle{14}
The determinant of the function $\zeta$ will be the product of all the differences between pairs of $(a+b)x-ab$. The number of these differences is
\[
\frac12m(m-1)(\tfrac12m(m-1)-1)=\frac14(m+1)m(m-1)(m-2)
\]
This number therefore gives the degree of the determinant of the function $\zeta$ with respect to the indeterminate $x$. The determinant of the function $z$ will have the same degree. The determinant of the function $Z$, however, may have a lower degree, whenever some coefficients, starting from the highest power of $x$, vanish. Our task is now to prove that \textit{all} the coefficients in the determinant of the function $Z$ certainly cannot vanish.

On examining more closely the differences whose product is the determinant of the function $\zeta$, we find that some of them---the differences between pairs of $(a+b)x-ab$ having an element in common---give
\[
\text{the product of all}\quad(a-b)(x-c)
\]
while the remaining ones---the differences between pairs of $(a+b)x-ab$ whose elements are distinct---give
\[
\text{the product of all}\quad(a+b-c-d)x-ab+cd,\quad\text{with repetitions excluded.}
\]
The first product clearly contains each factor $a-b$ exactly $m-2$ times, and each factor $x-c$ exactly $(m-1)(m-2)$ times. We therefore readily conclude that this product is
\[
=\pi^{m-2}\upsilon^{(m-1)(m-2)}
\]
If we denote the second product by $\rho$, the determinant of the function $\zeta$ is thus
\[
=\pi^{m-2}\upsilon^{(m-1)(m-2)}\rho
\]
Further, let $r$ denote the function of the indeterminates $x,l',l'',l'''$ etc. which becomes $\rho$ under the substitutions $l'=\lambda',l''=\lambda'',l'''=\lambda'''$ etc., and let $R$ denote
\SourceLeaf{056}
the function of $x$ alone into which $r$ passes under the substitutions $l'=L',l''=L''$, $l'''=L'''$ etc. It is then clear that the determinant of the function $z$ is
\[
=p^{m-2}y^{(m-1)(m-2)}r
\]
and that the determinant of the function $Z$ is
\[
=P^{m-2}Y^{(m-1)(m-2)}R
\]
Since, by hypothesis, $P$ is not $=0$, it now remains to prove that $R$ certainly cannot vanish identically.

\SecondArticle{15}
For this purpose we shall introduce one more indeterminate $w$ and consider the product of all
\[
(a+b-c-d)w+(a-c)(a-d)
\]
with repetitions excluded. Since it involves $a,b,c$ etc. symmetrically, it can be represented as a polynomial in the indeterminates $w,\lambda',\lambda'',\lambda'''$ etc. We shall denote this function by $f(w,\lambda',\lambda'',\lambda'''\text{ etc.})$. The number of the factors $(a+b-c-d)w+(a-c)(a-d)$ is
\[
=\frac12m(m-1)(m-2)(m-3)
\]
from which we readily conclude that
\[
f(0,\lambda',\lambda'',\lambda'''\text{ etc.})=\pi^{(m-2)(m-3)}
\]
and hence also
\[
f(0,l',l'',l'''\text{ etc.})=p^{(m-2)(m-3)}
\]
as well as
\[
f(0,L',L'',L'''\text{ etc.})=P^{(m-2)(m-3)}
\]
In general, the polynomial $f(w,\allowbreak L',\allowbreak L'',\allowbreak L'''\text{ etc.})$ will have degree
\[
\frac12m(m-1)(m-2)(m-3)
\]
but in special cases it may have a lower degree, if some coefficients, starting from the highest power of $w$, happen to vanish. It is impossible, however, for the entire function to be identically $=0$, since
\SourceLeaf{057}
the equation just found shows that at least its last term does not vanish. Suppose that the highest term of the polynomial $f(w,\allowbreak L',\allowbreak L'',\allowbreak L'''\text{ etc.})$ having a nonvanishing coefficient is $Nw^\nu$. If we substitute $w=x-a$, then $f(x-a,L',L'',L'''\text{ etc.})$ is clearly a polynomial in the indeterminates $x,a$, or, equivalently, a polynomial in $x$ with coefficients depending on the indeterminate $a$. Its highest term is nevertheless $Nx^\nu$, and thus has a fixed coefficient, independent of $a$, which is not $=0$. Likewise, $f(x-b,L',L'',L'''\text{ etc.}),f(x-c,L',L'',L'''\text{ etc.})$ are polynomials in the indeterminate $x$, each with highest term $Nx^\nu$, while the coefficients of the following terms depend respectively on $b,c$ etc. Hence the product of the $m$ factors
\[
\begin{gathered}
f(x-a,L',L'',L'''\text{ etc.})\\
f(x-b,L',L'',L'''\text{ etc.})\\
f(x-c,L',L'',L'''\text{ etc.})\\
\text{etc.}
\end{gathered}
\]
is a polynomial in $x$ whose highest term is $N^mx^{m\nu}$, while the coefficients of the following terms depend on the indeterminates $a,b,c$ etc.

Next consider the product of these $m$ factors:
\[
\begin{gathered}
f(x-a,l',l'',l'''\text{ etc.})\\
f(x-b,l',l'',l'''\text{ etc.})\\
f(x-c,l',l'',l'''\text{ etc.})\\
\text{etc.}
\end{gathered}
\]
As a function of the indeterminates $x,a,b,c$ etc., $l',l'',l'''$ etc., symmetric with respect to $a,b,c$ etc., it can be represented as a function of the indeterminates $x,\lambda',\lambda'',\lambda'''$ etc., $l',l'',l'''$ etc., which we shall denote by
\[
\varphi(x,\lambda',\lambda'',\lambda'''\text{ etc.},\ l',l'',l'''\text{ etc.})
\]
Thus
\[
\varphi(x,\lambda',\lambda'',\lambda'''\text{ etc.},\ \lambda',\lambda'',\lambda'''\text{ etc.})
\]
is the product of the factors
\[
\begin{gathered}
f(x-a,\lambda',\lambda'',\lambda'''\text{ etc.})\\
f(x-b,\lambda',\lambda'',\lambda'''\text{ etc.})\\
f(x-c,\lambda',\lambda'',\lambda'''\text{ etc.})\\
\text{etc.}
\end{gathered}
\]
\SourceLeaf{058}
and is therefore identically divisible by $\rho$, since it is easy to see that every factor of $\rho$ occurs in one of those factors. We therefore set
\[
\begin{aligned}
&\varphi(x,\lambda',\lambda'',\lambda'''\text{ etc.},\ \lambda',\lambda'',\lambda'''\text{ etc.})\\
&=\rho\psi(x,\lambda',\lambda'',\lambda'''\text{ etc.})
\end{aligned}
\]
where the symbol $\psi$ represents a polynomial. It follows readily that we also have identically
\[
\begin{aligned}
&\varphi(x,\allowbreak L',\allowbreak L'',\allowbreak L'''\text{ etc.},\allowbreak \ L',\allowbreak L'',\allowbreak L'''\text{ etc.})\\
&=R\psi(x,L',L'',L'''\text{ etc.})
\end{aligned}
\]
But we have proved above that the product of the factors
\[
\begin{gathered}
f(x-a,L',L'',L'''\text{ etc.})\\
f(x-b,L',L'',L'''\text{ etc.})\\
f(x-c,L',L'',L'''\text{ etc.})\\
\text{etc.}
\end{gathered}
\]
which is $=\varphi(x,\lambda',\lambda'',\lambda'''\text{ etc.},\ L',L'',L'''\text{ etc.})$, has highest term $N^mx^{m\nu}$. The function $\varphi(x,\allowbreak L',\allowbreak L'',\allowbreak L'''\text{ etc.},\allowbreak $
$L',L'',L'''\text{ etc.})$ therefore has the same highest term, and is certainly not identically $=0$. Hence $R$ cannot be identically $=0$ either, and neither can the determinant of the function $Z$. Q. E. D.

